\section{Results}
\label{sec:results}

\subsection{Primary Result 1: TempGlitch Pair-Disjoint Follow-up}

Table~\ref{tab:r5-bounded-results} reports the frozen pair-disjoint, non-locked, same-support
TempGlitch comparison. Every row uses 12 normal-negative and 22 buggy-positive evaluation
episodes after calibration on two normal episodes; source, pair, and episode overlap across roles
is zero.

\begin{table}[t]
\caption{Non-locked TempGlitch validation-only identical-episode family. Best observed rows within
the frozen R5 family.}
\label{tab:r5-bounded-results}
\centering
\small
\begin{tabularx}{\textwidth}{@{}lXccccX@{}}
\toprule
Family & Configuration & AUROC & AUPRC & F1 & FPR@95TPR & 95\% CI \\
\midrule
LeWM best AUROC & seed44 \texttt{lewm\_l2\_max} + episode \texttt{mean} & 0.696970 & 0.796125 & 0.723404 & 0.75 & AUROC [0.479152, 0.878319]; F1 [0.536585, 0.851984] \\
LeWM best AUPRC & seed44 \texttt{lewm\_l2\_top2\_mean} + episode \texttt{mean} & 0.670455 & 0.797118 & 0.711111 & 0.75 & AUROC [0.450536, 0.857792]; F1 [0.526316, 0.844487] \\
Baseline best AUROC & \texttt{feature\_distance} + episode \texttt{top2\_mean} & 0.617424 & 0.746156 & 0.086957 & 0.833333 & AUROC [0.397633, 0.810653]; F1 [0, 0.260870] \\
Baseline best AUPRC & \texttt{feature\_distance} + episode \texttt{mean} & 0.602273 & 0.777896 & 0.240000 & 1.00 & AUROC [0.407068, 0.800000]; F1 [0, 0.444444] \\
Baseline best F1 & \texttt{frame\_diff} + episode \texttt{mean} & 0.530303 & 0.701055 & 0.785714 & 0.833333 & AUROC [0.321329, 0.754941]; F1 [0.666667, 0.885246] \\
\bottomrule
\end{tabularx}
\end{table}


The best recorded LeWM-Glitch row is seed44 \texttt{lewm\_l2\_max} with mean episode aggregation.
It reaches AUROC 0.7159, AUPRC 0.8026, F1 0.7143, precision 0.7500, recall 0.6818, balanced
accuracy 0.6326, and FPR@95TPR 0.7500 with AUROC 95\% CI [0.5349, 0.8770]. The strongest simple
baseline by AUROC is \texttt{feature\_distance} with \texttt{top2\_mean} aggregation at
AUROC 0.6136 and AUPRC 0.7310. Within this exact frozen support, the best LeWM-Glitch row shows
stronger observed same-support separation than the simple baselines, but the AUROC intervals
overlap and the operating-point false-positive rate remains high.

\subsection{Exact-Support Learned Baselines (TempGlitch)}

Table~\ref{tab:k1-learned-baselines} extends the same TempGlitch follow-up support with three
learned video baselines.

\begin{table}[t]
\caption{K1 learned baselines on the frozen non-locked TempGlitch follow-up support. Every row
uses the same two calibration-normal episodes and the same 12 normal-negative / 22
buggy-positive evaluation episodes. Learned rows report the best aggregation per method over
\texttt{mean}, \texttt{max}, and \texttt{top2\_mean}.}
\label{tab:k1-learned-baselines}
\centering
\scriptsize
\begin{tabularx}{\textwidth}{@{}lXrrrrrrrX@{}}
\toprule
Family & Configuration & AUROC & AUPRC & F1 & Prec. & Recall & Bal. acc. & FPR@95 & 95\% CI \\
\midrule
LeWM & seed44 \texttt{lewm\_l2\_max} + episode \texttt{mean} & 0.7159 & 0.8026 & 0.7143 & 0.7500 & 0.6818 & 0.6326 & 0.7500 & AUROC [0.5349, 0.8770]; F1 [0.5854, 0.8293] \\
Simple baseline & \texttt{frame\_diff} + episode \texttt{mean} & 0.5833 & 0.7216 & 0.6154 & 0.7059 & 0.5455 & 0.5644 & 0.7500 & AUROC [0.4181, 0.7603]; F1 [0.4516, 0.7500] \\
Simple baseline & \texttt{feature\_distance} + episode \texttt{top2\_mean} & 0.6136 & 0.7310 & 0.1600 & 0.6667 & 0.0909 & 0.5038 & 0.8333 & AUROC [0.4636, 0.7545]; F1 [0, 0.3575] \\
Learned baseline & \texttt{video\_autoencoder} + episode \texttt{mean} & 0.5606 & 0.7005 & 0.6818 & 0.6818 & 0.6818 & 0.5492 & 0.8333 & AUROC [0.3909, 0.7386]; F1 [0.5556, 0.7907] \\
Learned baseline & \texttt{cnn\_lstm} + episode \texttt{max} & 0.6136 & 0.7249 & 0.7111 & 0.6957 & 0.7273 & 0.5720 & 0.9167 & AUROC [0.4405, 0.8000]; F1 [0.5946, 0.8163] \\
Learned baseline & \texttt{video\_transformer} + episode \texttt{mean} & 0.5909 & 0.7756 & 0.2400 & 1.0000 & 0.1364 & 0.5682 & 1.0000 & AUROC [0.4615, 0.7273]; F1 [0, 0.4286] \\
\bottomrule
\end{tabularx}
\end{table}


The strongest learned baseline is \texttt{cnn\_lstm} with \texttt{max} episode aggregation, at
AUROC 0.6136, AUPRC 0.7249, and F1 0.7111. Its AUROC numerically matches the strongest simple
baseline and remains below the best LeWM-Glitch row. The paired delta-AUROC versus the best
LeWM-Glitch row is $-0.1023$ with 95\% bootstrap interval $[-0.2425,\ 0.0372]$ and
DeLong $p=0.2091$, so the evidence supports only bounded stronger-observed-separation wording
for LeWM-Glitch rather than a significant superiority claim.

\subsection{Primary Result 2: R5-XGame Frozen Evaluation}

Table~\ref{tab:r5-xgame-results} reports the frozen non-locked R5-XGame family. This is a
co-primary headline result alongside TempGlitch, evaluated on a different game distribution
with a larger buggy-positive class (12 normal-negative, 60 buggy-positive episodes).

% R5-XGAME cross-dataset comparison table.
% SKELETON ONLY — do not replace placeholders until both R5-TempGlitch and
% R5-WOB outputs are validated and the R5-XGAME comparison script has run.
\begin{table}[t]
\centering
\caption{Cross-dataset comparison of episode-level glitch detection (R5-XGAME).
TempGlitch results are from the frozen non-locked R5 evaluation family.
World of Bugs results are pending validated R5-WOB output.
No cross-game generalization claim is supported until both datasets have
validated evidence.}
\label{tab:r5-xgame-results}
\begin{tabular}{llllll}
\toprule
Dataset & Method & Scorer & Aggregation & AUROC & AUPRC \\
\midrule
TempGlitch & LeWM (best) & TODO & TODO & TODO & TODO \\
TempGlitch & frame\_diff & --- & TODO & TODO & TODO \\
TempGlitch & feature\_distance & --- & TODO & TODO & TODO \\
\midrule
World of Bugs & LeWM (best) & TODO & TODO & TODO & TODO \\
World of Bugs & frame\_diff & --- & TODO & TODO & TODO \\
World of Bugs & feature\_distance & --- & TODO & TODO & TODO \\
\bottomrule
\end{tabular}
\end{table}


The best recorded LeWM-Glitch row (seed44 \texttt{lewm\_mse\_max} + \texttt{top2\_mean})
achieves AUROC 0.9097 and AUPRC 0.9814, with AUROC 95\% CI [0.8281, 0.9720]. Simple baselines
reach AUROC 0.7681 (\texttt{feature\_distance}) and 0.7167 (\texttt{frame\_difference}).
On this split, the best LeWM-Glitch AUROC is numerically above the best baseline AUROC, but the
AUROC confidence intervals still overlap ([0.8281, 0.9720] versus [0.6484, 0.8719]); this
supports bounded stronger observed separation, not significant superiority.

These values are not directly compared with TempGlitch values. The two result families differ
in dataset distribution, class balance, calibration support, action mode, and training lineage.
The R5-XGame result is evidence about its own frozen split; it does not establish cross-game
generalization.

\subsection{Bounded GlitchBench K2 Benchmark (Honest Negative)}

Table~\ref{tab:glitchbench-benchmark} reports the validated K2 GlitchBench slice. This is an
honest negative result that we report as a community-value finding.

\input{tables/glitchbench_benchmark}

On this image-level synthetic-normal split (24 validation examples),
\texttt{feature\_distance}, \texttt{video\_autoencoder}, \texttt{cnn\_lstm}, and
\texttt{video\_transformer} each reach AUROC 1.0 and AUPRC 1.0, while the best recorded
LeWM-Glitch row reaches AUROC 0.5. This is a clear negative result for the current LeWM
configuration on static image-level inputs: the latent prediction error signal does not provide
useful ranking on this split. At the same time, the non-LeWM AUROC-1.0 rows share F1 0.6667
and balanced accuracy 0.5 under the shared train-normal \texttt{p95} threshold rule,
highlighting the distinction between ranking quality and threshold calibration. K2 acts as a
bounded negative comparison surface that characterizes the current method's failure mode and
motivates future work on image-level LPE aggregation.

\subsection{K3 Mechanistic Ablation (Honest Negative)}

Table~\ref{tab:k3-ablation-results} summarizes the controlled K3 validation-normal readout.
This is an honest negative mechanistic finding reported as a community-value contribution.

\begin{table}[t]
\centering
\caption{Validated K3 mechanistic ablation readout on the frozen validation-normal support. Lower
prediction loss is better; SIGReg and action comparisons remain bounded mechanistic diagnostics,
not detection metrics.}
\label{tab:k3-ablation-results}
\small
\textbf{Panel A: group means}
\begin{tabular}{llrr}
\toprule
SIGReg & Action mode & Mean best total val. loss & Mean best prediction loss \\
\midrule
off & real & 0.001181 & 0.001181 \\
off & zero\_action & 0.001075 & 0.001075 \\
on & real & 0.280237 & 0.019721 \\
on & zero\_action & 0.284195 & 0.018489 \\
\bottomrule
\end{tabular}

\medskip
\textbf{Panel B: SIGReg-on minus SIGReg-off prediction-loss delta}
\begin{tabular}{lrr}
\toprule
Seed & Real-action delta & Zero-action delta \\
\midrule
42 & +0.016560 & +0.019388 \\
43 & +0.020677 & +0.017309 \\
44 & +0.018384 & +0.015546 \\
\bottomrule
\end{tabular}

\medskip
\textbf{Panel C: real-action minus zero-action prediction-loss delta}
\begin{tabular}{lrr}
\toprule
Seed & SIGReg-on delta & SIGReg-off delta \\
\midrule
42 & -0.003042 & -0.000215 \\
43 & +0.003361 & -0.000008 \\
44 & +0.003377 & +0.000539 \\
\bottomrule
\end{tabular}
\end{table}


Across the six matched SIGReg-on versus SIGReg-off comparisons, the lower-is-better prediction
loss never improves when SIGReg is enabled. Real-action versus zero-action prediction-loss
deltas are also not directionally stable across seeds. This is an honest negative mechanistic
readout for the specific K3 training objective; it does not invalidate the detection-performance
results above, but it does indicate that the SIGReg and action-conditioning components of
LeWM do not provide measurable benefits at this scale and configuration.

\subsection{Qualitative Surprise Timelines}

The qualitative timeline lane is summarized in Table~\ref{tab:qualitative-timelines}.

\begin{table}[t]
\centering
\caption{Qualitative-only P6 surprise timelines produced from the frozen non-locked TempGlitch
follow-up artifacts using seed44 \texttt{lewm\_l2\_max} with
\texttt{mean} aggregation. No temporal-localization metric is claimed.}
\label{tab:qualitative-timelines}
\scriptsize
\begin{tabularx}{\textwidth}{@{}l l X r r@{}}
\toprule
Label & Category & Source episode & Windows & Episode score \\
\midrule
Buggy & Shooting Error & Godot\_Shooting\_Error\_Buggy\_107 & 1088 & 1.696 \\
Buggy & Velocity Bug & Godot\_Teleportation\_TPS\_Buggy\_100 & 847 & 1.492 \\
Buggy & Frozen Animation & Godot\_Frozen\_Animation\_Buggy\_Sample\_13 & 861 & 1.485 \\
Normal & Shooting Error & Godot\_Shooting\_Error\_Normal\_101 & 1408 & 1.593 \\
Normal & Shooting Error & Godot\_Shooting\_Error\_Normal\_105 & 1401 & 1.228 \\
Normal & Shooting Error & Godot\_Shooting\_Error\_Normal\_107 & 1289 & 1.206 \\
\bottomrule
\end{tabularx}
\end{table}


These figures show that the best validated TempGlitch LeWM-Glitch configuration yields
interpretable within-episode score trajectories. However, the current artifacts do not expose
real temporal span labels. The timelines therefore remain qualitative illustrations only and do
not support temporal IoU, temporal AP, onset, offset, or duration claims. Temporal localization
is future work.
